\scanpage{31}
\mspage{\handwritten{- 46 -}}
% archive sheet number: 30

COROLLAIRE.- \uline{Soient} $\handwritten{\gamma}$ \uline{et} $\handwritten{\gamma}'$ \uline{deux routes achevées à valeurs dans} A, \uline{ayant même origine} $\handwritten{a}$ \uline{et même extrémité} $\handwritten{b}$, \uline{et telles qu'il existe une déformation de} $\handwritten{\gamma}$ \uline{en} $\handwritten{\gamma}'$ \uline{dans} A \uline{qui laisse fixes} $\handwritten{a}$ \uline{et} $\handwritten{b}$. \uline{Alors}, \uline{si l'expression différentielle} $\handwritten{\sum_k f}_k\,dx_k$ \uline{satisfait aux conditions} (1), \uline{on a} $\handwritten{\int_\gamma}(\handwritten{\sum_k f}_k\,dx_k)=\handwritten{\int_{\gamma'}}(\handwritten{\sum_k f}_k\,dx_k)$.

On peut supposer que $\handwritten{\gamma}$ et $\handwritten{\gamma}'$ sont définies dans le même intervalle $I=\handwritten{[\alpha,\beta]}$ ; soit $\handwritten{\gamma}^{\unsure{o|0|\circ}}$ l'opposé de $\handwritten{\gamma}$, et $\handwritten{\gamma}_1(t)=\handwritten{\gamma}_{\handwritten{\unsure{o|0|\circ}}}(t-\handwritten{\beta}+\handwritten{\alpha})$ un chemin équivalent à $\handwritten{\gamma}^{\handwritten{\unsure{o|0|\circ}}}$ défini dans $I+(\handwritten{\beta}-\handwritten{\alpha})$. Les routes $\handwritten{\gamma\vee\gamma}_1$ et $\handwritten{\gamma}'\handwritten{\vee\gamma}_1$ sont des circuits dans A et peuvent évidemment être déformés l'un en l'autre dans A en vertu des hypothèses ; en leur appliquant le th.1, on a aussitôt le corollaire.

\unsure{DÉFINITION|DÈFINITION} 5.- \uline{On dit qu'un domaine} $A\handwritten{\subset}K^n$ \uline{est contractile si tout lacet dans} A \uline{peut être déformé dans} A \uline{en un point de} A.

Par exemple, tout domaine \uline{étoilé} par rapport à un point est contractile.

Une expression différentielle $\handwritten{\sum_{k=1}^n f}_k\,dx_k$ à coefficients holomorphes dans un domaine A est dite \uline{fermée} si elle satisfait aux conditions (1). On voit donc que :

\unsure{THÈORÈME|THÉORÈME|THÈORÉME|THÉORÉME} 2 (Cauchy).- \uline{Dans un domaine contractile l'intégrale d'une expression différentielle fermée} (à coefficients holomorphes) \uline{le long d'un circuit quelconque est nulle}.

On peut évidemment, pour le cas des variables \uline{réelles}, généraliser le th.1 au cas où les $\handwritten{f}_k$ sont seulement supposées continument différentiables ; est-ce que cela en vaut la peine ?

\handwritten{\unsure{3|2}}. \uline{Primitives d'une fonction holomorphe}.

Soit $A\handwritten{\subset}K^n$ un domaine \uline{quelconque}, $\handwritten{g}$ une fonction holomorphe dans A ; si on pose $d\handwritten{g}=\handwritten{\sum_{k=1}^n f}_k\,dx_k$, il est immédiat que cette expression différentielle est \uline{fermée} ; en outre, en vertu des définitions, le long d'une route $\handwritten{\gamma}$ définie dans $I=\qquad$, on a $\handwritten{\int_\gamma}d\handwritten{g}=\handwritten{\int_\gamma}(\handwritten{\sum_k f}_k\,dx_k)=\handwritten{g}(\handwritten{\gamma}(\handwritten{\beta}))-\handwritten{g}(\handwritten{\gamma}(\handwritten{\alpha}\unsure{)|))|}$
% The source leaves the space after I= blank. At the right edge a faint arc may be one closing parenthesis; a second, or no intentional closing stroke, remains possible. Accents on DEFINITION and THEOREME are handwritten over the typed E letters.
